% HD-R13. GPT-6 Astra Pro /root; inherited effort not exposed.
% 21 November 2026. Actual products and a complete cubic-character analysis.
\section{HD-R13: external products of arbitrary threefold polynomial characters}
\label{sec:threefold-polynomial-products}

Let \(X_1,X_2\) be abelian threefolds, and let \(F_j\) be
actual perfect complexes with nonzero whole characters
\[
 \operatorname{ch}(F_j)=e^{B_j}P_j(H_j),\qquad
 B_j,H_j\in\operatorname{NS}(X_j)_{\mathbb Q},\qquad
 H_j^3\ne0,\quad P_j\in\mathbb Q[z],\quad\deg P_j\le3.
                                                               \tag{R13.1}
\]
Take their actual external product and any finite composition of
the actual Fourier--Mukai equivalences, isogeny pullbacks and
pushforwards, line twists, translations, shifts and derived duals
specified in HD-N17. Every resulting whole character
\[
             e^B\bigl(P(\theta)+\lambda t\bigr),\qquad
 0\ne\lambda t\in W_K(A),\qquad
 B\in\operatorname{NS}(A)_{\mathbb Q},                     \tag{R13.2}
\]
on a polarized Weil sixfold of signature \((3,3)\) has split
rational \(K\)-Hermitian form. The same conclusion holds for
a nonzero rational multiple of a transported character. This
extends HD-N17 beyond its two equal, prescribed imaginary-secant
factors: the polynomials, quadratic fields and nondegenerate
divisor classes of the two threefolds may differ. Pure and
double-root cubic cases are included below, with their actual
Clifford maps and Lie algebras. No existence of an additional
perfect complex is inferred from a formal character.

\paragraph{1. Explicit polynomial and Clifford changes.}
For \(V=H^1(X,\mathbb Q)\), use
\(\mathcal W=V\oplus V^*\),
\(Q(v+\xi,w+\eta)=\xi(w)+\eta(v)\), and
\(c(v+\xi)=v\wedge+\iota_\xi\). Multiplication by \(e^D\)
intertwines these operators through the literal rational map
\[
        v+\xi\longmapsto v-\iota_\xi D+\xi.                \tag{R13.3}
\]
This holds for every rational two-form \(D\). It is used as a
linear Clifford isometry when \(D\) is not the class of an
actual chosen line bundle; no such tensor functor is asserted.
Multiplication of a nonzero spinor by a nonzero rational scalar
leaves its vector and projective stabilizers and annihilator
unchanged, directly from their defining linear equations.
These are the precise maps used to remove \(e^{B_j}\) and
to compare polynomial coefficients.

First suppose the constant coefficient of \(P\) is zero.
Choose a rational Darboux basis for \(H\), so
\(H=\sum_{i=1}^3 a_i b_i\), and write
\(V_H=\int_X H^3/6\ne0\). In the dual basis the actual
Poincare kernel has class
\(\sum_i(a_i B_i-b_i A_i)\). Integrating its exponential
pair by pair gives
\[
 \int_{a_i b_i}e^{t a_i b_i+a_i B_i-b_i A_i}
             =t-A_iB_i,
 \qquad
 \mathrm{FM}^H(e^{tH})=V_H t^3e^{-\widehat H/t},
 \quad\widehat H=\sum_iA_iB_i.                            \tag{R13.4}
\]
The factor \(V_H\) accounts for the rational volume of the
Darboux basis. This is an identity of polynomials after truncation
in the threefold exterior algebra, so coefficient comparison gives
\[
 \mathrm{FM}^H(H^j/j!)
       =V_H\frac{(-\widehat H)^{3-j}}{(3-j)!}.
                                                               \tag{R13.5}
\]
The kernel is an actual Fourier--Mukai equivalence, with the
inverse and graph maps in Mukai, Theorem2.2 and formula(3.1),
at the precise locators of HD-N03. The formula also shows that
\(\widehat H\) is a rational algebraic nondegenerate divisor
class: it is a nonzero rational multiple of the degree-two
component of the actual kernel transform of \(H^2\).

Choose an integer \(a>0\) and a line bundle with class \(aH\).
The top-degree coefficient of \(e^{kaH}P(H)\) is a nonzero
polynomial in \(k\) of degree at most three. Among
\(k=0,1,2,3\), at least one makes it nonzero. The actual
line twist followed by this actual Fourier--Mukai equivalence
therefore gives a polynomial character in \(\widehat H\)
with nonzero constant coefficient. For a preexisting rational
shear \(B\), one may instead use the rational Clifford
composition (R13.3), followed by (R13.4); the proof below only
needs these linear intertwiners. Thus rank-zero polynomials are
covered without constructing a normalized object by division.

We may now compare the actual character, through the specified
rational scalar and shear maps, with
\[
                    w=1+uH^2/2+vH^3/6,
                        \qquad u,v\in\mathbb Q.             \tag{R13.6}
\]
Indeed dividing by its nonzero constant and applying
\(e^{-(\mu_1/\mu_0)H}\) makes the first factorial moment
zero. Every step is an explicit invertible map on the spin
module; its use in the obstruction proof does not assert an
object with that divided character.

\paragraph{2. The exact cubic cases.}
For \(u\ne0\), put \(a=v/u\) and
\(\Delta=a^2+4u\). When \(\Delta\ne0\), set
\[
 r_\pm=\frac{a\pm\sqrt\Delta}{2},\qquad
 \gamma_+=\frac{-r_-}{r_+-r_-},\qquad
 \gamma_-=\frac{r_+}{r_+-r_-}.
                                                               \tag{R13.7}
\]
Both coefficients are nonzero, since
\(\gamma_+\gamma_-=u/\Delta\ne0\). The four factorial
moments of
\(\gamma_+e^{r_+H}+\gamma_-e^{r_-H}\) are
\(1,0,u,v\): the first two are immediate, and the next
two follow from the exact recurrence \(r^2=ar+u\).
Consequently this is exactly \(w\), not only its leading
terms. Its two pure-spinor annihilators are the transverse
graphs \(\{-r_\pm H^\sharp\xi+\xi\}\).

For \(\Delta<0\), write \(s^2=\Delta/4=-d\),
\(d>0\), and remove the rational shear \(e^{aH/2}\).
The character becomes
\(\gamma_+e^{sH}+\overline{\gamma_+}e^{-sH}\) over the
imaginary quadratic field \(K_0=\mathbb Q(s)\).
The Fock map applying successive Clifford operators from the
second graph to the first exponential is exactly the map
in HD-N17, (N17.6), with \(C=H^\sharp\).
Its top exterior part is \(2sC\), hence invertible, so it
is an actual intertwiner. In that model the spinor is
\(b_0+b_6\mathrm{vol}_6\), both coefficients nonzero.
Degree two and degree four in its Lie action kill respectively
the wedge and contraction bivectors; degree zero and degree six
then force the middle trace and every projective scalar to be
zero. Thus its vector and projective stabilizers coincide.

The rational form and its explicit hyperbolic basis are
\[
 J=\begin{pmatrix}0&dC\\-C^{-1}&0\end{pmatrix},\qquad
 h_0(x,y)=Q(x,y)+s^{-1}Q(x,Jy),\qquad
                         \mathfrak s_w=\mathfrak{su}(h_0).
                                                               \tag{R13.8}
\]
Here \(J^2=-d\), \(s\) acts by \(J\), and
\(\{a_i,a_i^*:1\le i\le3\}\) is a \(K_0\)-basis
with hyperbolic Gram matrix. Directly,
\(Ja_i=b_i^*\), \(Ja_i^*=d b_i\); the remaining
basis vectors are obtained by applying \(J\).
Also \(h_0(a_i,a_j^*)=\delta_{ij}\), and all same-side
pairings vanish. This is the literal \(K_0\)-isometry
\(U^3\to(\mathcal W,h_0)\), extending the calculation of
HD-N17 without an ampleness or equal-polynomial restriction.

For \(\Delta>0\), the two graphs and the same Fock map
are defined over \(\mathbb R\). The identical degree
calculation gives real stabilizer \(\mathfrak{sl}_6(\mathbb R)\)
and zero projective scalar. For \(u=0,v\ne0\), (R13.6)
already has the rational vacuum/top form, so its stabilizer
is \(\mathfrak{sl}_6(\mathbb Q)\), again with zero
projective scalar. In all these stable cases the one-Clifford
map \(\mathcal W\to S^-\), \(z\mapsto c(z)w\), is
injective: in the relative Fock model its two components have
degrees one and five and are injective on the complementary
annihilator spaces.

For \(\Delta=0\), one has \(u=-r^2\ne0\),
\(v=-2r^3\), with \(r=a/2\in\mathbb Q^*\), and
\[
                         w=e^{rH}(1-rH).                    \tag{R13.9}
\]
After the displayed shear the character is \(1+D\) for
the rational nondegenerate two-form \(D=-rH\).
The remaining case \(u=v=0\) is the pure spinor \(1\).
These statements exhaust every nonzero polynomial after the
explicit changes in paragraph1.

\paragraph{3. The tangent case, including its projective stabilizer.}
On \(\bigwedge^*V\), write a spin Lie element as
\(\beta\wedge+D_A-\tfrac12\operatorname{tr}A+\iota_\gamma\).
For its action on \(1+D\) to equal \(c(1+D)\), its
degree-four component requires \(\beta\wedge D=0\).
The map \(D\wedge:\Lambda^2V\to\Lambda^4V\) is an
isomorphism by the following explicit matrix. Write
\(D=p_1+p_2+p_3\), \(p_i=a_i b_i\). Each of the twelve
two-monomials using different symplectic pairs is sent to
itself wedged with the unused \(p_i\); these images are
twelve distinct signed basis monomials. On the remaining
bases \((p_1,p_2,p_3)\) and
\((p_1p_2,p_1p_3,p_2p_3)\), its matrix is
\[
 \begin{pmatrix}1&1&0\\1&0&1\\0&1&1\end{pmatrix},
                         \qquad\det=-2.
\]
This gives its inverse over \(\mathbb Q\), so \(\beta=0\).
The degree-two and degree-zero components give exactly
\[
 D_AD=\ell D,\qquad \operatorname{tr}A=3\ell,
              c=-\ell/2,\qquad\iota_\gamma D=\ell.
                                                               \tag{R13.10}
\]
The trace identity follows by applying \(D_A\) to the
nonzero top volume \(D^3\). Equivalently,
\(A=(\ell/2)I+A_0\), with \(A_0\in\mathfrak{sp}(V,D)\).
The projective stabilizer has a surjective Lie map onto
\(\mathfrak{sp}_6(\mathbb Q)\), \(A\mapsto A_0\).
Its kernel is solvable: the primitive contraction bivectors
form an abelian ideal, and the remaining scalar parameter
\(\ell\) acts on them by scalar multiplication. The
vector stabilizer has \(\ell=0\) and is
\[
       \mathfrak{sp}_6(\mathbb Q)\ltimes\Lambda^2_0V^*.
                                                               \tag{R13.11}
\]
In particular its semisimple quotient has dimension21.
The one-Clifford map is injective also here: for
\(c(v+\xi)(1+D)=0\), the degree-three part gives
\(v\wedge D=0\), whence \(v=0\), and the degree-one
part gives \(\iota_\xi D=0\), whence \(\xi=0\).

\paragraph{4. The product stabilizer and the target annihilator.}
For two nonpure factors, the spin Lie algebra on their
orthogonal Clifford sum has the actual decomposition
\[
 \mathfrak{so}(\mathcal W_1)\oplus
 \mathfrak{so}(\mathcal W_2)\oplus
                         (\mathcal W_1\otimes\mathcal W_2).
                                                               \tag{R13.12}
\]
The mixed term acts on \(w_1\otimes w_2\) by the tensor
of the two injective one-Clifford maps. It is therefore
injective into \(S_1^-\otimes S_2^-\); the other terms
land in \(S_1^+\otimes S_2^+\). Hence a stabilizing
element has no mixed term. The remaining equation forces
each factor action to be scalar on its own spinor, and the
two scalars sum to zero. This gives the exact product
stabilizer as the fiber product of the two explicitly
computed projective stabilizers over their scalar characters.
Its solvable radical has semisimple quotient formed from
the stable dimension35 factors and the tangent
\(\mathfrak{sp}_6\) factors just computed.

A pure factor is separately excluded by an actual annihilator
calculation on the target. Complexify its \(K\)-action and
write \(H^1(A,\mathbb C)=V_+\oplus V_-\), each of
dimension six. The polarization has weight zero for the
linear action with weights \(+1,-1\) on these summands,
and \(t\) is a sum of two nonzero top volumes of weights
\(+6,-6\). The two coefficients are both nonzero because
the original nonzero class is rational and conjugation exchanges
the volumes. For
\(c(v+\xi)(P(\theta)+\lambda t)=0\), its degree-five
components of weights \(+5,-5\) are precisely
\(\lambda\iota_\xi t\). All polynomial contributions
in that degree have weights \(+1,-1\). Contraction with
each six-volume is injective on its six-dimensional dual,
so \(\xi=0\). The degree-seven components of weights
\(+5,-5\) are then \(\lambda v\wedge t\). Wedge with
the opposite volume is injective on each six-space, so
\(v=0\). Thus the full target Clifford annihilator is
zero. The shear \(e^B\) preserves this fact through
(R13.3). A product with a pure factor has the nonzero
annihilator of that factor inside its first Clifford summand.
Every actual transport used here intertwines the Clifford
action, or its derived-dual Lie and parity version as in
HD-N17. For the latter, the reversal operator
\(\tau|_{\Lambda^jV}=(-1)^{j(j-1)/2}\) satisfies on
every even spinor \(w\) the literal identity
\(c(v+\xi)\tau w=\tau c(v-\xi)w\).
Thus the invertible map \((v,\xi)\mapsto(v,-\xi)\)
also preserves annihilator dimension under derived dual.
Consequently such a product cannot reach (R13.2).

\paragraph{5. The target Lie algebra and the actual hyperbolic isometry.}
For the nonpure cases, \(\mathfrak{su}(h)\) of the target
is a rational simple Lie algebra of dimension35. Its actual
contragredient action on \(H^1(A)\) preserves \(\theta\)
and both determinant lines, so it kills (R13.2), after
conjugation by the explicit shear. The scalar-extension,
simplicity and faithful-action maps are those proved in full
in HD-N17, paragraph6. All actual transports identify the
stabilizers by (N17.11), the explicit isogeny map, line
shear and derived-dual Lie automorphism in its paragraph5.
Their hypotheses concern these actual functors and are unchanged
by the differing polynomials here.

The resulting injection into the product stabilizer cannot have
image in its solvable radical: a nonzero simple Lie algebra
is perfect and nonsolvable, and every subalgebra of a solvable
Lie algebra is solvable. Its projection to some semisimple
factor is therefore nonzero, hence injective. A tangent
\(\mathfrak{sp}_6\) factor has dimension21 and cannot receive
that injection. A stable real factor is excluded by the
explicit real Killing signatures
\[
 \mathfrak{su}(3,3):(18,17),\qquad
 \mathfrak{sl}_6(\mathbb R):(20,15).                       \tag{R13.13}
\]
The first follows from its18-dimensional noncompact off-diagonal
part and17-dimensional compact part; the second from trace-zero
symmetric matrices of dimension20 and skew matrices of dimension15.
An injective map between these dimension35 real algebras would
be an isomorphism and preserve its Killing form, contradicting
these explicit signatures.

The nonzero projection must consequently be a rational Lie
isomorphism to one of the imaginary-secant factors
\(\mathfrak{su}(U^3)\) in (R13.8), possibly over a different
imaginary quadratic field from the target. Take the preimages
of its three commuting hyperbolic diagonal matrices with entries
\(+1,-1\) on the respective hyperbolic pairs. Their adjoint
eigenvalues are integral. The exact argument in HD-N17,
paragraph7, applies without an equality of the two quadratic
fields: each preimage is semisimple, its standard eigenvalue
differences are rational, and trace zero makes its individual
eigenvalues rational. The three preimages are therefore
simultaneously diagonalizable over the target \(K\).
Their joint weights span a three-dimensional rational space,
and skew Hermitianity pairs only opposite weights. In target
dimension six there are exactly three opposite one-dimensional
pairs. Choosing vectors \(v_i,w_i\) in these spaces and
replacing \(w_i\) by \(w_i/h(v_i,w_i)\) gives the
actual Gram matrix \(U^3\). Sending its standard basis to
these vectors is the explicit target \(K\)-linear Hermitian
isometry. This proves the stated split conclusion.

The operation on whole characters is a computation in the
exterior realization of coherent \(K\)-theory, with every
actual kernel, pullback, norm and tensor functor retained.
It closes these external-product suppliers for the remaining
nonsplit sixfold condition-A problem (cote16 scan050,
SIGN\_NEXT SN-C03/C04). New coupled differentials and
other whole characters retain their own construction gate.
